\SPBGContentSlide
  {Методы и обозначения}
  {%
    \centering
    \SmallBody
    \renewcommand{\arraystretch}{1.08}
    \begin{tabularx}{.95\linewidth}{@{}l>{\raggedright\arraybackslash}X>{\raggedright\arraybackslash}X@{}}
      \toprule
      Код & Метод & Смысл шага \\
      \midrule
      H1 & Базовый MDM & перенос веса между двумя точками \\
      H2 & SMO & парное изменение двойственных коэффициентов \\
      H3 & MDM-G & перенос в группу минимумов \\
      H4 & MDM-A & ускорение по повторяющимся циклам \\
      H5 & Треугольный MDM & выбор ближайшей точки локального треугольника \\
      H6 & Комбинированный MDM & выбор лучшего из нескольких кандидатов \\
      \bottomrule
    \end{tabularx}
  }

\SPBGContentSlide
  {Методы и обозначения}
  {%
    \centering
    \SmallBody
    \renewcommand{\arraystretch}{1.15}
    \begin{tabularx}{.95\linewidth}{@{}>{\raggedright\arraybackslash}X>{\raggedright\arraybackslash}X>{\raggedright\arraybackslash}X@{}}
      \toprule
      Обозначения & Группа методов & Постановка \\
      \midrule
      MDM, MDM-G, MDM-A & мягкие MDM-варианты & ограничения $0\leq u_j\leq C$ \\
      SMO, SVC, LinearSVC & SVM-решатели & мягкая линейная постановка \\
      \bottomrule
    \end{tabularx}
  }
